\section{A Local-First Kanban Board}
\label{sec:a-local-first-kanban-board}

The TypeScript library developed for this thesis lets a web application accept edits on one device and share them with others when a connection is available.
A Kanban board serves as its demonstration application.
Each card represents a ticket: its column shows how far the work has progressed, and its priority indicates how urgently the team should address it.
Moving a card changes its status; opening it lets a team member edit its details, assign it to someone, or leave a comment.

Alice opens the board on her laptop and Bob opens it on his phone.
Both see the ticket \enquote{Fix the sign-in error} in the \emph{Backlog} column, as shown in Figure~\ref{fig:kanban-board}.
Its priority is \emph{Medium}.

\begin{illustration}[htbp]
	\centering
	\caption{The ticket before Alice takes her laptop offline}
	\label{fig:kanban-board}
	\includegraphics[width=\linewidth]{kanban-board.png}
	\figuresource{Author's screenshot of the example application. The equals sign on the card denotes medium priority.}
\end{illustration}
\FloatBarrier

Alice loses her connection just as she starts investigating the error.
She drags the card to \emph{In progress}, and her board shows it in the new column immediately.
Meanwhile, Bob learns that the error prevents several people from signing in.
On his phone the ticket is still in \emph{Backlog}.
He opens it and raises its priority to \emph{High}.
At this point each person can see their own edit, but neither can see the other's.

\begin{table}[htbp]
	\centering
	\caption{Status and priority of the same ticket on each device}
	\label{tab:kanban-ticket}
	\begin{tabular}{@{}lll@{}}
		\toprule
		& Alice's laptop & Bob's phone \\
		\midrule
		Before disconnection & Backlog, Medium & Backlog, Medium \\
		After both edits, laptop offline & In progress, Medium & Backlog, High \\
		After synchronization & In progress, High & In progress, High \\
		\bottomrule
	\end{tabular}
\end{table}
\FloatBarrier

When Alice reconnects, Bob sees the card move and Alice sees its priority rise.
Table~\ref{tab:kanban-ticket} records the three states observed in two separate browser sessions running the application.
The outcome differs from the concurrent renames in Section~\ref{sec:conflict-free-replicated-data-types}: a ticket can display only one title, but Alice changed its status and Bob changed its priority, so both edits can take effect.
How does the library record a move on a disconnected laptop so that another device can incorporate it later?
